\documentclass[aps,prl,twocolumn,superscriptaddress,hyperref]{revtex4-2}
\usepackage{amsmath,amssymb,amsfonts}

\begin{document}

\title{Information-Capacity Bound on Non-Markovian Backflow}

\author{Huang Zhongchang}
\affiliation{Independent Researcher, Nanjing, China}

\begin{abstract}
For a bipartite quantum system whose pointer basis $\{|0\rangle,|1\rangle\}$ is einselected by a pure-dephasing interaction, we prove a combinatorial upper bound on non-Markovian information backflow: $P_{\rm reflux} \leq \min(q_S/q_E,\,(1-q_E)/q_E)$, where $q_S = \langle 0|\bar{\rho}_S|0\rangle$ and $q_E = \langle 0|\bar{\rho}_E|0\rangle$ are per-qubit pointer-basis populations before irreversible transfer. Forward transfer is a Lindblad jump $L = |1\rangle\langle 1|_S \otimes |1\rangle\langle 0|_E$ that only creates $|1\rangle$ states; the irreversibility of $|0\rangle \to |1\rangle$ follows directly from the jump operator, not from $H_{\rm int}$. The bound uses only counting. No physical constants enter. It connects finite-capacity constraints to established non-Markovianity measures in a complementary role, and is in principle falsifiable.
\end{abstract}

\maketitle

%%%%%%%%%%%%%%%%%%%%%%%%%%%%%%%%%%%%%%%%%%%%%%%%%%%%%%%%%%%%%%%%%%%%%%
\section{Introduction}
%%%%%%%%%%%%%%%%%%%%%%%%%%%%%%%%%%%%%%%%%%%%%%%%%%%%%%%%%%%%%%%%%%%%%%

While general relativity follows from the equivalence principle and special relativity from the constancy of $c$, quantum decoherence has resisted reduction to a single generative constraint. Non-Markovianity---the backflow of information from environment to system---is quantified by the Breuer-Laine-Piilo (BLP) trace-distance integral~\cite{Breuer:2009} and classified by the Buscemi et al.\ causal/noncausal revival distinction~\cite{Buscemi:2025}. BLP integrates the positive derivative of the trace distance to detect whether backflow occurred. Buscemi uses process-tensor tomography to determine whether a revival is causal. Both are phenomenological. Neither supplies an upper bound grounded in the information capacity of the physical substrate itself.

The question this Letter asks is whether there exists a limit on backflow that follows not from a system's coupling to its environment, but from the finite information capacity of its constituent degrees of freedom. The answer is yes, and the bound takes a surprisingly simple form.

We prove $P_{\rm reflux} \leq \min(q_S/q_E,\,(1-q_E)/q_E)$: the reflux fraction cannot exceed a ratio of pointer-basis populations. The proof is four lines of counting and requires no physical constants, no free parameters, and no assumptions about spectral densities or coupling strengths. What it does require is the working hypothesis that each qubit encodes at most 1 bit of distinguishable information in its pointer basis. The bound is not derivable from standard open-systems theory because standard theory does not track per-element information capacity.

%%%%%%%%%%%%%%%%%%%%%%%%%%%%%%%%%%%%%%%%%%%%%%%%%%%%%%%%%%%%%%%%%%%%%%
\section{Model}
%%%%%%%%%%%%%%%%%%%%%%%%%%%%%%%%%%%%%%%%%%%%%%%%%%%%%%%%%%%%%%%%%%%%%%

\subsection{Binary information capacity}

It is not obvious that nature admits a description in terms of binary information units. But it is also not obvious that such a description is impossible. We postulate that each qubit encodes 1 bit of distinguishable information in its pointer basis: the distinction between $|0\rangle$ and $|1\rangle$. This is the minimal nontrivial information capacity---less than 1 bit would mean no distinguishable states, and a degree of freedom with zero distinguishability cannot participate in physical dynamics~\cite{Shannon:1948}. Whether nature satisfies this postulate is an empirical question.

\subsection{Einselection of the pointer basis}

A pointer basis is not a mathematical convenience. In any system with a dominant pure-dephasing interaction
\begin{equation}
H_{\rm int} = \sum_{i} \sigma_z^{(i)} \otimes B_i,
\label{eq:hint}
\end{equation}
where $\sigma_z^{(i)} = |0\rangle\langle 0| - |1\rangle\langle 1|$, the basis $\{|0\rangle,|1\rangle\}$ is physically selected by environment-induced superselection~\cite{Zurek:2003}. Equation~\eqref{eq:hint} preserves all pointer-basis populations identically. Its role is solely to pick the basis; it contributes no population transfer. Population dynamics are governed by the separate Lindblad channel introduced in Sec.~\ref{sec:forward}.

From Section~\ref{sec:model} we carry forward two properties: the 1-bit hypothesis and the einselected pointer basis. The third piece---the dynamics of information transfer---is specified next.

\subsection{Forward transfer}
\label{sec:forward}

Forward transfer is the process by which a determined qubit influences an undetermined neighbor. It is modeled by the Lindblad jump operator
\begin{equation}
L_{S\to E} = |1\rangle\langle 1|_S \otimes |1\rangle\langle 0|_E,
\label{eq:jump}
\end{equation}
acting on a specific $S$--$E$ pair. When the source is $|1\rangle$ and the target is $|0\rangle$, the jump toggles the target to $|1\rangle$; otherwise $L_{S\to E}$ acts as the null operator. The source qubit is unchanged. The jump operator encodes exactly 1 bit of change.

Equation~\eqref{eq:jump} is a modeling choice. Its key property is that it only creates $|1\rangle$ states. Once a qubit transitions $|0\rangle \to |1\rangle$, no subsequent dynamics can return it to $|0\rangle$, because no term in either $H_{\rm int}$ or $L$ contains a $|0\rangle\langle 1|$ factor. The irreversibility is a property of the full GKSL dynamics: with $H_{\rm int}$ preserving pointer-basis populations and $L$ only creating $|1\rangle$, no term in the master equation generates $|1\rangle \to |0\rangle$ transitions. This is a strong condition satisfied when determination corresponds to an irreversible record-formation event. In systems with non-negligible $|1\rangle \to |0\rangle$ relaxation, the bound becomes approximate with corrections controlled by the $T_1$ time.

%%%%%%%%%%%%%%%%%%%%%%%%%%%%%%%%%%%%%%%%%%%%%%%%%%%%%%%%%%%%%%%%%%%%%%
\section{Theorem: Reflux Bound}
\label{sec:model}
%%%%%%%%%%%%%%%%%%%%%%%%%%%%%%%%%%%%%%%%%%%%%%%%%%%%%%%%%%%%%%%%%%%%%%

\subsection{Setup}

Partition the qubits into a system $S$ ($N_S$ qubits) and an environment $E$ ($N_E$ qubits). Let $q_S = N_S^{-1}\sum_{i\in S} \langle 0|\rho_i|0\rangle$ and $q_E$ similarly be pre-transfer pointer-basis populations. Section~\ref{sec:model} established two properties: (i) the pointer basis $\{|0\rangle,|1\rangle\}$ is the binary description of each qubit, and (ii) the transition $|0\rangle \to |1\rangle$ occurs at most once per qubit. The theorem that follows uses only these two properties.

Forward transfer: determined $S$-qubits toggle undetermined $E$-qubits via $L_{S\to E}$. Backflow: determined $E$-qubits toggle undetermined $S$-qubits via $L_{E\to S}$.

\subsection{Statement and proof}

\begin{quote}
\textbf{Theorem (Reflux Bound).} After irreversible forward transfer has occurred, the reflux fraction satisfies $P_{\rm reflux} \leq \min(q_S/q_E,\,(1-q_E)/q_E)$ when $N_S = N_E$.
\end{quote}

\textbf{Proof.}
The forward transfer capacity is the number of $E$-qubits available to receive determination: $C_F \equiv N_E \, q_E$. Backflow sources are determined $E$-qubits, numbering $N_E(1-q_E)$. Backflow targets are undetermined $S$-qubits, numbering $N_S q_S$. Each can be toggled at most once, so $N_{\rm back} \leq \min(N_S q_S,\; N_E(1-q_E))$. Defining $P_{\rm reflux} \equiv N_{\rm back} / C_F$,
\begin{equation}
P_{\rm reflux} \leq \frac{\min(N_S q_S,\; N_E(1-q_E))}{N_E q_E}
= \min\!\left(\frac{q_S}{q_E},\; \frac{1-q_E}{q_E}\right),
\label{eq:proof}
\end{equation}
where the last equality uses $N_S = N_E$.

$\square$

The inequality is the pigeonhole principle. The receiver-side constraint $q_S/q_E$ limits backflow by the availability of undetermined $S$-qubits. The source-side constraint $(1-q_E)/q_E$ limits it by the availability of determined $E$-qubits to serve as backflow sources. Neither constraint is redundant.

\subsection{Domain}

The theorem requires $C_F > 0$ ($q_E > 0$) and irreversible forward transfer. When $q_E = 1$, all $E$-qubits are undetermined and the premise is not met. The bound constrains the total reflux fraction, not individual events.

%%%%%%%%%%%%%%%%%%%%%%%%%%%%%%%%%%%%%%%%%%%%%%%%%%%%%%%%%%%%%%%%%%%%%%
\section{Relation to Non-Markovianity Measures}
%%%%%%%%%%%%%%%%%%%%%%%%%%%%%%%%%%%%%%%%%%%%%%%%%%%%%%%%%%%%%%%%%%%%%%

The BLP measure~\cite{Breuer:2009, Breuer:2016} integrates the positive derivative of the trace distance. It detects whether backflow occurred. S1 bounds how much can occur, given the finite capacity of the constituent qubits. The bound also differs from intrinsic-decoherence models~\cite{Milburn:1991}, which postulate a universal minimum dephasing rate from discrete time evolution. S1 constrains a specific mechanism---determination-mediated backflow---rather than imposing a global decoherence floor.

Buscemi et al.~\cite{Buscemi:2025} classify revivals as causal or noncausal. Neither framework supplies a capacity-based upper bound on the backflow fraction. The three approaches are complementary: detection, classification, and capacity-limited amplitude.

%%%%%%%%%%%%%%%%%%%%%%%%%%%%%%%%%%%%%%%%%%%%%%%%%%%%%%%%%%%%%%%%%%%%%%
\section{Discussion}
%%%%%%%%%%%%%%%%%%%%%%%%%%%%%%%%%%%%%%%%%%%%%%%%%%%%%%%%%%%%%%%%%%%%%%

We do not claim that nature is fundamentally composed of 1-bit elements. We claim that if each qubit encodes at most 1 bit of distinguishable information in its pointer basis, then the reflux fraction is bounded as above. The combinatorial argument is classical---it uses only binary distinguishability and the pigeonhole principle. What makes the bound relevant to quantum systems is that the pointer basis, the irreversibility of determination, and the per-qubit population $q$ are all specified by the quantum mechanical structure of the system-environment interaction. Without einselection there is no preferred binary basis; without $L$ there is no irreversible toggle; without both, the bound has no domain of applicability. Quantum mechanics supplies the physical conditions under which a classical counting bound constrains a real system.

The bound is non-trivial---providing a constraint tighter than $P_{\rm reflux} \leq 1$---only when both $q_S < q_E$ and $q_E > 1/2$. Outside this regime, at least one of the two constraints exceeds unity and reduces to the trivial bound. This is not a defect of the inequality but a statement about when finite capacity actually limits backflow: when the environment has substantial remaining capacity ($q_E > 1/2$) and the system is more determined than the environment ($q_S < q_E$).

Several limitations are explicit. The bound constrains total reflux, not per-event backflow. It assumes $N_S = N_E$ for the simplified form. When $q_E < 1-q_E$, the source constraint $(1-q_E)/q_E$ can be tighter than $q_S/q_E$. The $L$ operator is a modeling choice whose validity depends on the dominance of the determination mechanism over competing relaxation channels. The pure-dephasing form of $H_{\rm int}$ is an idealization; non-dephasing corrections produce $O(\varepsilon^2\tau^2)$ deviations.

%%%%%%%%%%%%%%%%%%%%%%%%%%%%%%%%%%%%%%%%%%%%%%%%%%%%%%%%%%%%%%%%%%%%%%
\section{Conclusion}
%%%%%%%%%%%%%%%%%%%%%%%%%%%%%%%%%%%%%%%%%%%%%%%%%%%%%%%%%%%%%%%%%%%%%%

Given the 1-bit hypothesis and a pure-dephasing pointer-basis interaction, the reflux fraction satisfies $P_{\rm reflux} \leq \min(q_S/q_E,\,(1-q_E)/q_E)$. The bound is parameter-free, requires no physical constants, and connects finite information capacity to the established theory of non-Markovianity. Its experimental test awaits.

\begin{thebibliography}{10}

\bibitem{Breuer:2009}
H.-P.~Breuer, E.-M.~Laine, and J.~Piilo,
Measure for the degree of non-Markovian behavior,
Phys.\ Rev.\ Lett.\ \textbf{103}, 210401 (2009).

\bibitem{Breuer:2016}
H.-P.~Breuer, E.-M.~Laine, J.~Piilo, and B.~Vacchini,
Colloquium: Non-Markovian dynamics in open quantum systems,
Rev.\ Mod.\ Phys.\ \textbf{88}, 021002 (2016).

\bibitem{Buscemi:2025}
F.~Buscemi, R.~Gangwar, K.~Goswami, H.~Badhani, T.~Pandit, B.~Mohan, S.~Das, and M.~N.~Bera,
Causal and noncausal revivals of information,
PRX Quantum \textbf{6}, 020316 (2025).

\bibitem{Zurek:2003}
W.~H.~Zurek,
Decoherence, einselection, and the quantum origins of the classical,
Rev.\ Mod.\ Phys.\ \textbf{75}, 715 (2003).

\bibitem{Shannon:1948}
C.~E.~Shannon,
A mathematical theory of communication,
Bell Syst.\ Tech.\ J.\ \textbf{27}, 379 (1948).

\bibitem{Milburn:1991}
G.~J.~Milburn,
Intrinsic decoherence in quantum mechanics,
Phys.\ Rev.\ A \textbf{44}, 5401 (1991).

\bibitem{SM}
See Supplemental Material for the complete proof, physical origin of the Lindblad jump operator, and detailed comparison with BLP and Buscemi measures.

\end{thebibliography}

\end{document}
